% Generated draft: requires mathematical and source review.
\documentclass[11pt]{article}
\usepackage[margin=25mm]{geometry}
\usepackage{fontspec}
\setmainfont{Latin Modern Roman}
\usepackage{amsmath,amssymb,mathtools,unicode-math}
\setmathfont{Latin Modern Math}
\usepackage{enumitem,needspace,xurl,hyperref}
\setlength{\parindent}{0pt}
\setlength{\parskip}{5pt}
\setlength{\emergencystretch}{3em}
\newcommand{\N}{\mathbb{N}}
\newcommand{\Z}{\mathbb{Z}}
\newcommand{\Q}{\mathbb{Q}}
\newcommand{\R}{\mathbb{R}}
\newcommand{\C}{\mathbb{C}}
\newcommand{\F}{\mathbb{F}}
\newcommand{\A}{\mathbb{A}}
\newcommand{\PP}{\mathbb{P}}
\newcommand{\E}{\mathbb{E}}
\newcommand{\eps}{\varepsilon}
\newcommand{\dd}{\,\mathrm{d}}
\begin{document}
\section*{UnsolvedMath ID: 800704}
\textbf{Source problem code:} catalogue entry
\subsection{Definitions and mathematical statement}
For a finitely generated field extension \(K/F\), call \(K/F\) \emph{rational} if \(K\cong_F F(t_1,\ldots,t_r)\) for some integer \(r\geq 0\), where the \(t_i\) are algebraically independent over \(F\). Call it \emph{stably rational} if there exist integers \(r,s\geq 0\) and algebraically independent indeterminates such that
\[
K(u_1,\ldots,u_r)\cong_F F(v_1,\ldots,v_s).
\]
Call it \emph{retract rational} if there exist a finitely generated integral \(F\)-subalgebra \(R\subseteq K\) with \(\operatorname{Frac}(R)=K\), an integer \(s\geq 0\), a nonzero polynomial \(f\in F[t_1,\ldots,t_s]\), and unital \(F\)-algebra homomorphisms
\[
R\xrightarrow{\alpha}F[t_1,\ldots,t_s][1/f]\xrightarrow{\beta}R,
\qquad \beta\circ\alpha=\operatorname{id}_R.
\]
Two extensions \(K/F\) and \(L/F\) are \emph{stably \(F\)-isomorphic} if \(K(u_1,\ldots,u_r)\cong_F L(v_1,\ldots,v_s)\) for some finite sets of algebraically independent indeterminates.

The supplied source asks whether each of the extensions
\[
Z(F,n)/F,\qquad Z_m(F,n)/F,\qquad Z_H(F,n)/F
\]
is rational, stably rational, or retract rational, with the three properties considered separately. However, the supplied record does not define these fields, specify the hypotheses on \(F\), or give the admissible values and meaning of the parameters \(n,m,H\). Thus this is a conditional statement of the intended rationality questions: the target extensions and their parameter ranges must be recovered from the cited source before a fully specified quantified problem can be stated. The definitions above do not supply those missing constructions or assert that the displayed extensions satisfy their prerequisites.

As background, the source reports a result attributed to Katsylo and Schofield: for admissible indices satisfying \(n=ab\) and \(\gcd(a,b)=1\), the field \(Z(F,n)\) is stably isomorphic to the compositum \(Z(F,a)Z(F,b)\). The supplied record does not specify the embeddings or ambient field used for this compositum, nor the full hypotheses of that result. It also reports that similar statements are possible for \(Z_m(F,n)\) and \(Z_H(F,n)\), without stating them precisely. The source qualifies the resulting reduction to prime-power indices as applying often, but not always.
\subsection{Short English statement}
The problem asks whether the source's fields \(Z(F,n)\), \(Z_m(F,n)\), and \(Z_H(F,n)\) are rational, stably rational, or retract rational over \(F\). Their constructions and parameter restrictions are absent from the supplied record, so the intended scope remains unresolved. The registry labels the problem open; that label is supplied metadata, not an independently verified status assessment.
\subsection{Sources}
\begin{itemize}
\item Supplied local UnsolvedMath registry record: permanent ID \(800704\), exact source problem code \texttt{catalogue entry}, title \emph{Open Problems on Central Simple Algebras --- Problem 14}. The record attributes the entry to Auel, Brussel, Garibaldi, and Vishne, with proposed year 2010, and labels it \texttt{open}.
\item Auel, Brussel, Garibaldi, and Vishne, \emph{Open Problems on Central Simple Algebras} (2010), Problem 14. Supplied source URL: \url{https://arxiv.org/abs/1006.3304}. The provenance records access on 2026-10-16 and extraction method \texttt{arxiv\_latex}; these are supplied provenance claims, not independently verified access or source checks.
\item The question carries the supplied reference label \texttt{procesi:67}, p.~240. No full bibliographic entry for this label is supplied.
\item The background attributes the relatively prime factorization result to Katsylo, reference label \texttt{Katsylo}, and Schofield, reference label \texttt{Schofield}. Full bibliographic entries and precise theorem hypotheses are not supplied.
\end{itemize}
\end{document}
